\chapter{Taylor coefficients and changes of order}
\label{chap:ordered-taylor}

A mixed coefficient of a polynomial is one scalar. An operation with
ordered inputs must distribute that number among several words.
For the mixed cubic, the ordered contraction placed the coefficient
of \(x_1x_2^2\) on the word \(212\), while the words \(122\) and
\(221\) gave zero. The same rule holds for every potential: sort
the indices, then exchange the first two. The other operations,
whose inputs include products of primitive generators, retain
additional information.

The choice of order also enters the diagonal operator itself.
Exchanging two adjacent variables changes its divided differences.
A finite conjugation implements that change. Two sequences of
exchanges can give different conjugations, even when they induce
the same permutation. Their difference is removed by an explicit
homotopy of differential-algebra maps. Thus order has both a
visible effect on formulas and a controlled effect on comparisons.

\section{An ordered contraction in any number of variables}
\label{sec:ot-contraction}

Let \(k\) be a commutative algebra over \(\mathbb Q\). Fix an
integer \(d\ge1\), and put
\[
 R=k[y_1,\ldots,y_d],\qquad S=R[r_1,\ldots,r_d],
 \qquad x_i=y_i+r_i.
\]
The formal version uses \(R=k[[y_1,\ldots,y_d]]\) and
\(S=R[[r_1,\ldots,r_d]]\). Take a polynomial \(W\) in the
polynomial case and a formal power series \(W\) in the formal case.
No regularity assumption on its gradient is imposed.
All degrees in this chapter are parities. Derivatives in \(y\)
keep the coordinates \(x\) fixed.

On the exterior algebra of \(d\) odd vectors let \(\beta_i\)
denote left exterior multiplication and let \(\alpha_i\) denote
contraction. Contraction deletes the \(i\)-th vector with the
sign of its position, and gives zero if it is absent. Hence
\[
 \{\alpha_i,\alpha_j\}=\{\beta_i,\beta_j\}=0,\qquad
 \{\alpha_i,\beta_j\}=\delta_{ij},
\]
where \(\{A,B\}=AB+BA\) for odd operators.
Define
\[
 g_i(x,y)=
 \frac{W(y_1,\ldots,y_{i-1},x_i,\ldots,x_d)
       -W(y_1,\ldots,y_i,x_{i+1},\ldots,x_d)}
      {x_i-y_i}.
\]
The numerator is divisible by \(x_i-y_i\), with a unique quotient:
multiplication by the formal or polynomial variable \(r_i\)
shifts coefficients and is injective even if \(k\) has zero divisors.
Telescoping gives \(\sum_i r_i g_i=W(x)-W(y)\). Therefore
\[
 D=\sum_i(r_i\alpha_i+g_i\beta_i),\qquad
 D^2=(W(x)-W(y))1.
\]
The algebra \(E=\operatorname{End}_S(S\otimes_k\Lambda(k^d))\)
has differential
\(\delta A=DA-(-1)^{|A|}AD\). Its square is zero because
\(D^2\) is scalar.

Put \(H_i=-\partial_{y_i}D\). Differentiating the square of \(D\)
gives
\begin{gather}
 \delta H_i=W_i(y),\qquad \delta\beta_i=r_i,\qquad
 \{H_i,\beta_j\}=\delta_{ij},\label{eq:ot-primitives}\\
 \{H_i,H_j\}=c_{ij},\qquad
 c_{ij}=-W_{ij}(y)-\sum_l r_l\,\partial_{y_i}\partial_{y_j}g_l.
 \label{eq:ot-clifford}
\end{gather}
Here \(W_i=\partial_iW\) and \(W_{ij}=\partial_i\partial_jW\).
For an increasing subset \(I\) let \(H_I\) be the ordered product
of its \(H_i\), and define \(\beta_J\) in the same way.
The monomials \(H_I\beta_J\) form an \(S\)-basis of \(E\).
Indeed, the ordered products \(\alpha_I\beta_J\) are a basis,
as follows by tensoring the four matrix units for one exterior
variable. Replacing \(\alpha_i\) by
\(H_i=\alpha_i-\sum_j(\partial_{y_i}g_j)\beta_j\)
changes this basis triangularly, with diagonal entries one.

Let \(C\) be the free \(R\)-module on symbols \(\xi_I\), with
\(\xi_\varnothing=1\), parity \(|I|\), and differential
\[
 d_C\xi_I=\sum_{a=1}^{|I|}(-1)^{a-1}W_{i_a}(y)
                 \xi_{I\setminus\{i_a\}}.
\]
The scalar values in \eqref{eq:ot-primitives} identify the matrix
complex with this gradient complex tensored with the Koszul
complex on the normal coordinates \(r_i\).
Define \(i(\xi_I)=H_I\). Define \(p(fH_I\beta_J)\) to be
\(f(y,y)\xi_I\) if \(J\) is empty, and zero otherwise.
Both are chain maps and \(pi=1_C\).

For \(f\in S\), write \(f^{(j)}\) for the result of setting
\(x_l=y_l\) for all \(l<j\), and put
\[
 \Delta_j f=\frac{f^{(j)}-f^{(j)}|_{x_j=y_j}}{r_j}.
\]
With \(\min\varnothing=d+1\), set
\begin{equation}\label{eq:ot-homotopy}
 h(fH_I\beta_J)=(-1)^{|I|}
       \sum_{j<\min J}\Delta_j f\,H_I\beta_j\beta_J.
\end{equation}
This formula has
\[
 \delta h+h\delta=1-ip,\qquad ph=hi=h^2=0.
\]
For one normal coordinate the homotopy is
\(\beta(f(r)-f(0))/r\). Its anticommutator with the differential
is evaluation error \(f-f(0)\) on functions and is the identity
on terms containing \(\beta\). For several coordinates let \(P_j\)
evaluate \(r_j=0\) and kill its exterior generator. The tensor
homotopy is \(h_1+P_1h_2+\cdots+P_1\cdots P_{d-1}h_d\),
with the Koszul signs for odd maps. Its anticommutator telescopes
to \(1-P_1\cdots P_d\), giving \eqref{eq:ot-homotopy}.
Each squared summand vanishes. Every mixed product contains
\(h_jP_j=0\) or \(P_jh_j=0\), which proves \(h^2=0\).
The other side identities follow from evaluation.

Let \(\mathsf s\) denote suspension. Use the tensor signs and
deconcatenation from Section~\ref{sec:cmc-transfer}, and define
\[
 F_1=\mathsf s i\mathsf s^{-1},\qquad b_1=\mathsf s d_C\mathsf s^{-1},
 \qquad
 \mathcal R_n=\sum_{j=1}^{n-1}B_2(F_j\otimes F_{n-j}),
\]
\[
 F_n=-\mathsf s h\mathsf s^{-1}\mathcal R_n,\qquad
 b_n=\mathsf s p\mathsf s^{-1}\mathcal R_n,\qquad
 B_2(\mathsf s A,\mathsf s B)=(-1)^{|A|}\mathsf s(AB).
\]
The contraction identities are precisely those used in the
proof in Section~\ref{sec:cmc-transfer}. That proof applies
without a dimension restriction: the relative-coderivation defect
vanishes by induction on length, and the split injection
\((\mathsf s i\mathsf s^{-1})^{\otimes n}\) then gives \(b^2=0\).
It also proves strict unitality and the morphism equation.
Thus \(b\) is an \(A_\infty\)-structure and \(F:C\to E\)
is an \(R\)-linear \(A_\infty\)-quasi-isomorphism.
In the formal case each fixed-arity component is continuous:
products, evaluation, and divided differences are continuous,
and only finitely many trees occur in each arity.
Word length remains a direct sum, with no convergence assertion
across arities.

\section{The order carrying a Taylor coefficient}
\label{sec:ot-coefficients}

The triangular dependence of the divided differences controls
the ordered words. The coefficient \(g_j\) uses only
\(x_j,\ldots,x_d\) and \(y_1,\ldots,y_j\).
Consequently
\(\partial_y^\nu g_j=0\) if \(\nu\) uses an index greater
than \(j\). Its other derivatives are free of
\(x_1,\ldots,x_{j-1}\).
By \eqref{eq:ot-homotopy}, both \(p\) and \(h\) kill
\[
 fH_i\beta_j
 \quad\text{if \(f\) is free of \(x_1,\ldots,x_{j-1}\),}
\]
and they kill \(f\beta_j\beta_l\) if \(f\) is free of
all \(x_a\) with \(a<\min(j,l)\).
These vanishing statements remove all but one split of a
primitive word in the transfer recurrence.

For a multi-index \(\nu=(\nu_1,\ldots,\nu_d)\), put
\[
 |\nu|=\sum_i\nu_i,\qquad \nu!=\prod_i\nu_i!,\qquad
 W_\nu(y)=\frac{\partial^\nu W(y)}{\nu!},\qquad
 D_\nu=\frac{\partial_y^\nu D}{\nu!}.
\]
When \(|\nu|\ge2\), the \(\alpha_i\)-part of \(D\) has
vanished under differentiation, so
\(D_\nu=\sum_j(\partial_y^\nu g_j/\nu!)\beta_j\).

\begin{lemma}\label{lem:ot-remainder}
Suppose \(|\nu|\ge2\), \(i\ge\max\operatorname{supp}\nu\),
and \(\varphi=\partial_y^\nu g_i/\nu!\).
Then
\[
 \varphi(y,y)=W_{\nu+e_i}(y),\qquad
 h(\varphi\,1)=D_{\nu+e_i},
\]
where \(e_i\) is the multi-index with a single \(1\) in position \(i\).
\end{lemma}

\begin{proof}
For a series \(G\) in one variable define its Taylor remainder
\[
 T_mG(x,y)=\sum_{j\ge m}\frac{G^{(j)}(y)}{j!}(x-y)^{j-m}.
\]
Expansion of monomials proves Taylor's formula for polynomials;
coefficientwise passage proves it for formal series.
Removing the first term and differentiating at fixed \(x\) give
\[
 \frac{T_mG-T_mG|_{x=y}}{x-y}=T_{m+1}G,\qquad
 \partial_y T_mG=mT_{m+1}G.
\]
In the derivative, the coefficient of
\(G^{(j)}(y)(x-y)^{j-m-1}\) is
\(1/(j-1)!-(j-m)/j!=m/j!\).

Write \(\nu=\nu'+ae_i\), where \(\nu'\) uses indices below \(i\).
Set
\[
 G(t)=\frac{1}{\nu'!}\partial_{y}^{\nu'}
             W(y_1,\ldots,y_{i-1},t,x_{i+1},\ldots,x_d).
\]
The divided difference defining \(g_i\) is \(T_1\) in the
pair \((x_i,y_i)\). The preceding derivative identity gives
\(\varphi=T_{a+1}G\). Evaluation at \(x_i=y_i\), then at
the other diagonal coordinates, yields \(W_{\nu+e_i}\).

The coefficient of \(\beta_j\) in \(h(\varphi)\) is
\(\Delta_j\varphi\). It is zero for \(j<i\), as is the
coefficient of \(\beta_j\) in \(D_{\nu+e_i}\).
For \(j=i\), the remainder identity gives
\(\Delta_i\varphi=T_{a+2}G
 =\partial_y^{\nu+e_i}g_i/(\nu+e_i)!\).
For \(j>i\), the earlier diagonal substitutions in
\(\varphi^{(j)}\) give
\(\partial^{\nu+e_i}W(y_{<j},x_{\ge j})/(\nu+e_i)!\).
Taking its divided difference in \(x_j\) gives the coefficient
of \(\beta_j\) in \(D_{\nu+e_i}\), since those derivatives
act on earlier coordinates and commute with this divided difference.
\end{proof}

For \(|\nu|=n\ge2\), list its indices in nondecreasing order
\(m_1,\ldots,m_n\), and define
\[
 w(\nu)=(m_2,m_1,m_3,\ldots,m_n).
\]
Repeated indices are retained. Write \(a_i=\mathsf s\xi_i\)
and \(e=\mathsf s1\). The \(a_i\) are even and \(e\) is odd.

\begin{theorem}\label{thm:ot-ordered-taylor}
The binary multiplication on \(C\) is the associative product
determined by ordered monomials and the relations
\[
 \xi_i\star\xi_j+\xi_j\star\xi_i=-W_{ij}(y).
\]
In particular \(\xi_i\star\xi_i=-W_{ii}(y)/2\).
For every primitive word \((i_1,\ldots,i_n)\) with multiplicities
\(\nu\),
\[
 F_n(a_{i_1},\ldots,a_{i_n})=
 \begin{cases}
 -\mathsf sD_\nu,&(i_1,\ldots,i_n)=w(\nu),\\
 0,&\text{otherwise}
 \end{cases}
 \qquad(n\ge2),
\]
and
\[
 b_n(a_{i_1},\ldots,a_{i_n})=
 \begin{cases}
 W_\nu(y)e,&(i_1,\ldots,i_n)=w(\nu),\\
 0,&\text{otherwise}
 \end{cases}
 \qquad(n\ge3).
\]
The differential \(b_1\) and all operations on composite inputs
remain those of the full contraction.
\end{theorem}

\begin{proof}
Reduce products \(H_IH_J\) using \eqref{eq:ot-clifford}.
Projection evaluates the central coefficients \(c_{ij}\) at
\(r=0\), where they are \(-W_{ij}(y)\). Thus the projected
product is the displayed Clifford product in its ordered basis.
Its associativity follows by reducing triple products and
evaluating the resulting polynomial identities.
The gradient differential is a derivation of this product:
it preserves every defining relation, since the \(W_i(y)\)
are central scalars.

On two primitive inputs \(\mathcal R_2(a_i,a_j)=-\mathsf sH_iH_j\).
For \(i<j\) the product is an ordered representative, killed by \(h\).
For \(i<j\), reversing the indices introduces \(c_{ij}\), and
\[
 h(c_{ij})=-\sum_l\partial_{y_i}\partial_{y_j}g_l\,\beta_l
          =-D_{e_i+e_j}.
\]
To check the first equality, apply \(\Delta_l\) to
\eqref{eq:ot-clifford}. Earlier terms contain a normal coordinate
already set to zero; later terms are free of \(x_l\).
The \(l\)-th term gives its coefficient.
For equal indices the same calculation includes the factor
\(1/2\), giving \(-D_{2e_i}\). This proves the formula for \(F_2\).

Suppose the formula for \(F\) holds in smaller arities.
A split with both blocks of length at least two gives a product
of \(\beta\)-linear expressions. Their coefficients are free of
the earlier \(x\)'s, so both \(p\) and \(h\) kill it.
A split with a singleton on the left gives \(H_jD_\mu\),
also killed by both maps. The only remaining split isolates
the last letter \(i\), after a word \(w(\mu)\).
It contributes \(\mathsf sD_\mu H_i\). If
\(D_\mu=\sum_j\varphi_j\beta_j\), then
\[
 D_\mu H_i=\varphi_i-\sum_j\varphi_jH_i\beta_j.
\]
The sum is killed by \(p,h\). If
\(i<\max\operatorname{supp}\mu\), its scalar term is zero.
Otherwise Lemma~\ref{lem:ot-remainder} gives the displayed
values of \(b_n,F_n\).
Appending a largest index to \(w(\mu)\) gives exactly
\(w(\mu+e_i)\). This proves the assertion by induction.
\end{proof}

For \(W=x_1^2x_2\), the coefficient \(W_{(2,1)}=1\)
therefore occurs in \(b_3(a_1,a_1,a_2)=e\), while the
other two orderings vanish. For \(W=x_1x_2x_3\), the
only nonzero primitive ternary value is
\(b_3(a_2,a_1,a_3)=e\). The rule concerns this contraction.
The cyclic change of coordinates in the preceding chapter
changes how coefficients are distributed among words.

If \(\nabla W(0)=0\), coefficientwise specialization at \(y=0\)
gives a minimal model of \(\operatorname{End}(D(x,0))\).
This follows by specializing the displayed contraction itself:
all its identities are \(R\)-linear and survive specialization.
It does not require preservation of arbitrary quasi-isomorphisms
by a nonflat quotient.

\begin{corollary}\label{cor:ot-quadratic}
If \(W\) has degree at most two, every operation \(b_n\)
with \(n\ge3\) vanishes on all inputs, including composite
inputs. The inclusion \(i\) is a strict differential-algebra
quasi-isomorphism from the differential Clifford algebra \(C\)
to \(E\).
\end{corollary}

\begin{proof}
Each \(g_i\) has degree at most one. The coefficients
\(\partial_{y_i}g_j\) and \(c_{ij}\) are constant, so
the ordered representatives \(H_I\) are closed under
multiplication with coefficients in \(R\).
The homotopy kills every such product. Hence \(F_2=0\),
and the recurrence gives \(F_n=0\) for all \(n\ge2\) and
\(b_n=0\) for all \(n\ge3\).
The contraction proves that \(i\) is a quasi-isomorphism.
The differential \(d_C\xi_i=W_i(y)\) remains part of the
model and need not vanish.
\end{proof}

\section{Exchanging two variables}
\label{sec:ot-exchange}

The telescoping formula can use any order of the variables.
Suppose \(i,j\) are adjacent in that order. Hold earlier
coordinates at \(y\) and later coordinates at \(x\), and denote
the resulting two-variable series by \(V(s,t)\).
Its mixed divided difference is
\[
 q_{ij}=
 \frac{V(x_i,x_j)-V(y_i,x_j)-V(x_i,y_j)+V(y_i,y_j)}
      {r_i r_j}.
\]
Each division is coefficientwise and has a polynomial or formal
quotient. Let \(D\) denote the diagonal with \(i\) before \(j\),
and let \(D'\) denote the diagonal with \(j\) before \(i\).
Subtracting their divided differences gives
\[
 g'_i-g_i=-r_jq_{ij},\qquad
 g'_j-g_j=r_iq_{ij}.
\]
All other coefficients agree. Put
\[
 U=1-q_{ij}\beta_i\beta_j,\qquad U^{-1}=1+q_{ij}\beta_i\beta_j.
\]
These are inverse even matrices since \((\beta_i\beta_j)^2=0\).
The Clifford relations give
\[
 [\beta_i\beta_j,D]=r_j\beta_i-r_i\beta_j,\qquad
 [\beta_i\beta_j,[\beta_i\beta_j,D]]=0.
\]
Consequently
\begin{equation}\label{eq:ot-conjugation}
 D'=UDU^{-1}.
\end{equation}
Conjugation \(A\mapsto UAU^{-1}\) is therefore a unital
\(S\)-linear isomorphism of differential algebras
\((E,\delta_D)\to(E,\delta_{D'})\). It preserves parity
and is continuous in the formal case.

This comparison retains the primitive maps as well.
For any retained coordinate \(y_l\), let
\(A_l=(\partial_{y_l}U)U^{-1}\), an even matrix.
Differentiating \eqref{eq:ot-conjugation} at fixed \(x\) yields
\[
 -\partial_{y_l}D'
   =U(-\partial_{y_l}D)U^{-1}+\delta_{D'}A_l.
\]
The primitive changes by a displayed boundary, including its
coefficient dependence.

There is also an explicit comparison on the relative Hochschild
cochains
\[
 C_R^n(E,E)=\operatorname{Hom}_{R,\mathrm{cont}}
                     (E^{\widehat\otimes_R n},E).
\]
The hat denotes completion in the normal variables of the tensor
factors; it is omitted in the polynomial case.
For a cochain \(f\) of arity \(n\), define
\[
 f'(A'_1,\ldots,A'_n)
  =U f(U^{-1}A'_1U,\ldots,U^{-1}A'_nU)U^{-1}.
\]
For arity zero the formula is conjugation of the output.
The differential, cup product, and every brace insertion commute
with this map: substitute the formula into each operation and
cancel adjacent \(U^{-1}U\). Since \(U\) is even, this cancellation
does not change any Koszul sign. The inverse uses \(U^{-1}\).
The statement holds for relative continuous cochains with the
coefficient ring and fixed-arity completions used above.

\section{Two exchange paths and their homotopy}
\label{sec:ot-paths}

Already three variables distinguish equality from compatibility
up to homotopy. For \(W=x_1x_2x_3\), reverse the order by
either sequence
\[
 123\longrightarrow213\longrightarrow231\longrightarrow321,
 \qquad
 123\longrightarrow132\longrightarrow312\longrightarrow321.
\]
Multiplying the corresponding matrices gives \(U_A=1+B_A\)
and \(U_B=1+B_B\), where
\[
 B_A=-x_3\beta_1\beta_2-y_2\beta_1\beta_3-x_1\beta_2\beta_3,
 \qquad
 B_B=-y_3\beta_1\beta_2-x_2\beta_1\beta_3-y_1\beta_2\beta_3.
\]
Products of two such exterior quadratic expressions vanish
in three variables. Thus
\[
 B_A-B_B=-r_3\beta_1\beta_2+r_2\beta_1\beta_3-r_1\beta_2\beta_3
        =-\delta_D(\beta_1\beta_2\beta_3).
\]
The two matrices are generally different. Their difference is
controlled by an exterior cubic expression.

To state the general comparison, a polynomial homotopy between
differential-algebra maps \(f_0,f_1:A\to B\) means a
differential-algebra map
\[
 f:A\longrightarrow k[\lambda,d\lambda]\otimes_k B
\]
whose evaluations at \(\lambda=0,1\), with \(d\lambda=0\),
are \(f_0,f_1\). Here \(\lambda\) is even, \(d\lambda\) is odd,
\((d\lambda)^2=0\), and the differential sends \(\lambda\)
to \(d\lambda\). The tensor product uses the Koszul sign rule.
In the formal case the normal and retained coefficients are completed,
while dependence on \(\lambda\) is polynomial.

\begin{theorem}\label{thm:ot-path-homotopy}
Two finite sequences of adjacent exchanges with the same initial
and final orders induce polynomially homotopic isomorphisms of
the diagonal endomorphism algebras.
The homotopy is \(S\)-linear and continuous in the formal case.
\end{theorem}

\begin{proof}
All exterior quadratic expressions commute with each other.
Write the product of the conjugating matrices along a path as
\(U=\exp B\), where \(B\) is the sum of their exterior
quadratic exponents. The exponential is a finite sum:
an exterior expression of degree greater than \(d\) vanishes.
Furthermore \([B,[B,D]]=0\), since \([B,D]\) is
exterior-linear. Thus the final operator is \(D+[B,D]\).

Let the two paths give \(B_0,B_1\) and put \(Z=B_1-B_0\).
Equality of the final operators gives \(\delta_D Z=[D,Z]=0\).
On expressions using only the \(\beta_i\), the differential
is the Koszul differential \(\delta_D\beta_i=r_i\).
Apply its contraction \eqref{eq:ot-homotopy}, with \(I\)
empty, and put \(Y=hZ\). Projection kills \(Z\), so
\[
 \delta_D Y=Z.
\]
The matrix \(Y\) is odd and exterior-cubic.
It is polynomial in the polynomial case and continuous in
the normal coordinates in the formal case.

Let \(D'\) be the common final operator, and set
\[
 U(\lambda)=\exp(B_0)\exp(\lambda Z),\qquad
 C_\lambda(A)=U(\lambda)AU(\lambda)^{-1}.
\]
Every \(C_\lambda\) is a differential-algebra isomorphism,
since \(Z\) commutes with \(D\).
For a homogeneous \(A\) define
\[
 K_\lambda(A)=C_\lambda\bigl(YA-(-1)^{|A|}AY\bigr).
\]
This map is odd and is a derivation along \(C_\lambda\):
\[
 K_\lambda(AB)=K_\lambda(A)C_\lambda(B)
                 +(-1)^{|A|}C_\lambda(A)K_\lambda(B).
\]
The identity \(\delta_DY=Z\), followed by the graded
commutator rule, gives
\[
 \partial_\lambda C_\lambda
  =C_\lambda[Z,-]
  =\delta_{D'}K_\lambda+K_\lambda\delta_D.
\]
Both identities can also be obtained by expanding the two products;
the terms containing \(Y\delta_D A\) cancel with opposite signs.

Define \(f(A)=C_\lambda(A)+d\lambda\,K_\lambda(A)\),
writing \(d\lambda\) on the left. The derivation identity proves
multiplicativity, because moving \(d\lambda\) past
\(C_\lambda(A)\) introduces \((-1)^{|A|}\).
For the differential,
\[
 d f(A)=C_\lambda(\delta_D A)
       +d\lambda\bigl(\partial_\lambda C_\lambda(A)
                            -\delta_{D'}K_\lambda(A)\bigr)
       =f(\delta_D A).
\]
The map is unital since \(K_\lambda(1)=0\).
Its endpoint evaluations are the two conjugations.
All exponentials are finite, and the only divided differences
occur in \(hZ\), so the asserted coefficient and continuity
properties follow.
\end{proof}

In the three-variable example, choose the second path as
endpoint zero and the first as endpoint one. Then
\(Z=-\delta_D(\beta_1\beta_2\beta_3)\) and
\(Y=-\beta_1\beta_2\beta_3\).
This gives the entire polynomial homotopy by substitution.
The cubic term is required because the two quadratic exponents
are not equal.

The same calculation has a formula for every potential.
Suppose \(i,j,k\) are three consecutive variables, and hold
the other coordinates at the values prescribed by the initial
order. For the resulting three-variable series \(V\), define
the mixed divided difference
\[
 q_{ijk}=
 \frac{\displaystyle
   \sum_{\epsilon\in\{0,1\}^3}
    (-1)^{3-\epsilon_1-\epsilon_2-\epsilon_3}
    V(z_i^{\epsilon_1},z_j^{\epsilon_2},z_k^{\epsilon_3})}
 {r_i r_j r_k},
 \qquad z_l^1=x_l,\quad z_l^0=y_l.
\]
The numerator is divisible by each normal coordinate.

\begin{corollary}\label{cor:ot-braid}
For the exchange sequences
\[
 ijk\to jik\to jki\to kji,\qquad
 ijk\to ikj\to kij\to kji,
\]
let \(B_A,B_B\) be the quadratic exponents of their
conjugating matrices. Then
\[
 B_A-B_B=-q_{ijk}\,
             \delta_D(\beta_i\beta_j\beta_k).
\]
The polynomial homotopy of Theorem~\ref{thm:ot-path-homotopy}
can therefore use \(Y=-q_{ijk}\beta_i\beta_j\beta_k\).
\end{corollary}

\begin{proof}
The coefficient of \(\beta_i\beta_j\) in the difference
is \(-q_{ij}(x_k)+q_{ij}(y_k)=-r_kq_{ijk}\).
The coefficient of \(\beta_i\beta_k\) is
\(-q_{ik}(y_j)+q_{ik}(x_j)=r_jq_{ijk}\).
The coefficient of \(\beta_j\beta_k\) is
\(-q_{jk}(x_i)+q_{jk}(y_i)=-r_iq_{ijk}\).
These are exactly the coefficients of
\(-q_{ijk}\delta_D(\beta_i\beta_j\beta_k)\).
Since \(\delta_D\) is \(S\)-linear, it kills \(q_{ijk}\),
and the asserted \(Y\) has differential \(B_A-B_B\).
\end{proof}

Each ordered transferred model maps by its explicit
\(A_\infty\)-quasi-isomorphism to its diagonal endomorphism
algebra. The conjugations therefore give actual zigzags between
the models associated to different orders. The theorem also
compares the differential-algebra maps obtained from different
exchange paths. A compatible family of higher homotopies among
these path homotopies requires further equations. The closed
formula for arbitrary composite inputs and the transport of a
specified cyclic cocycle are separate questions.
